\section{Experiments, predictive realization and stopping}\label{sec:experiments}
\subsection{Raw objects and executable tests}
All history, preparation, hidden-state and report spaces are standard Borel. The finite nonempty action set $A$ is the same at equivalent histories; any action restriction or public resource record belongs to the interface being compared. Raw prepared kernels are positive and normalized. Actions use only already available history and independent private randomization. No program reads an unknown physical parameter. The quantitative theorem below concerns a known instrument and known preparation; the morphism theorem separately quantifies over parameter families.

Fix legal finite future experiments and bounded terminal readouts, closed under their legal one-step continuations. A finite separating affine test realization means the following. The constant test and a finite vector $u(h)$ of executable-test expectations are Borel, their attained image has a Borel realization section, and equality of these coordinates is equivalent to equality of all stipulated future tests. Realizable random mixtures form a compact convex base $B\subset\R^d$, full-dimensional in this test chart. For every action the joint unnormalized expectation of each coordinate test and the next report is affine in the incoming coordinates. These are hypotheses to be verified, not consequences of positivity alone.

Write $C=\{w(1,u):w\ge0,u\in B\}$ and $e(w,v)=w$. The affine closure is represented by integrable positive linear instruments $L_a(y)$ on this finite cone, relative to reference probabilities $\lambda_a$; $\int L_a\,d\lambda_a$ preserves $e$. Define
\begin{equation}\label{eq:instrument}
 L_a(y)(1,u)=(\ell_a(y,u),U_a(y,u)),\qquad
 F_a(y,u)=U_a(y,u)/\ell_a(y,u).
\end{equation}
All raw identities hold off one reference-null set for each action, simultaneously for every $u$. If the likelihood is zero, positivity and bounded coordinates force $U_a(y,u)=0$; that event contributes zero without division. The actual acquired row and its operator are
\begin{equation}\label{eq:row}
 \mu_{a,u}=(F_a(\cdot,u))_\#(\ell_a(\cdot,u)\lambda_a),
 \qquad K_af(u)=\int f(v)\mu_{a,u}(dv).
\end{equation}
The realized test state is not supplied as a costless register to a finite observer.

\begin{lemma}[Measurable predictive realization]\label{lem:quotient}
The attained test image realizes the predictive quotient. It is minimal among sufficient Borel statistics whose attained images have Borel realization sections. The report and update kernels descend jointly in the finite action variable.
\end{lemma}
\begin{proof}
The separating hypothesis identifies precisely the equivalence classes. Compose a realization section of another sufficient statistic with $h\mapsto u(h)$; sufficiency makes the value independent of the section representative. This gives Borel factorization. Formula~\eqref{eq:instrument}, with the common null sets, gives the descended kernels. Successive integration realizes all legal finite continuations. This proves minimality in the stated category; it does not claim that an arbitrary measurable quotient is standard Borel, or that $B$ is an acquired support.
\end{proof}

\subsection{One protected audit}
A fixed executable audit $Y^{\rm aud}$ has conditional mean $u$ and affine second moment $c(u)$. A response $r\in B$, committed before the audit, has loss
\begin{equation}\label{eq:loss}
 \E[\norm{Y^{\rm aud}-r}^2\mid u]=c(u)-2\langle u,r\rangle+\norm r^2,
 \qquad g(u)=c(u)-\norm u^2.
\end{equation}
Assume $0\le\norm{Y^{\rm aud}-r}^2\le L_0$. The actual audit noise remains in $c$. For a categorical audit one uses its affine probability hyperplane with the induced Euclidean metric. For quantum experiments only this final audit is executed; no simultaneous intermediate measurement of incompatible tests is assumed.

Fix $s\in(0,1]$ and put $\beta=1-s$. After each acquisition, and after the finite update but before the audit, an independent fresh coin terminates with probability $s$. Thus
\begin{equation}\label{eq:stop}
 \Pp(\tau=t)=s\beta^{t-1},\quad t\ge1,\qquad \E\tau=s^{-1}.
\end{equation}
The hazard is independent of the physical state, reports and observer randomizations, conditional on every pre-coin history. Until the coin terminates, it conveys only that the run continues. The process is stopped and audited once. An unaudited reference execution is only a coupling device for these finite stopped laws. There is no state-dependent selection of favorable stopping reports.

For a full-history strategy $\pi$, let $X_t$ be its exact predictive state and set
\begin{equation}\label{eq:occupation}
 \nu_s^\pi=s\sum_{t\ge1}\beta^{t-1}\Law_\pi(X_t),\quad
 B_s^\pi=\int g\,d\nu_s^\pi,\quad V_s(u_0)=\inf_\pi B_s^\pi.
\end{equation}
These are probability laws and risks of the actual stopped experiment. The controller may adaptively choose each acquisition action, but cannot choose the imposed hazard.

A checkpoint receives the full history at the request and maps it to at most $M$ labels, with one label-only decoder. Its earlier acquisition workspace is not counted as $M$. A homogeneous observer has at most $M$ total labels, an initial label, fixed action probabilities $p_s$, one Borel report-update row at each label, and a fixed readout for each label. Only that label is retained between calls. An external stop request reads it. Its infimum risk is $\calR_{\rm hom}(s,M;u_0)$.

For the internally stopped architecture, $\calR_{\rm int}(s,M;u_0)$, halt status is part of the label. At each report the same prescribed independent coin selects halt versus continuation; no external observer clock is supplied. Halt labels have no further acquisition action. Section~\ref{sec:stop-strata} charges their separate copies. The complete call count in either core architecture is $\tau+2$ including initial preparation and final audit; hence its mean is $s^{-1}+2$ and it has no finite deterministic cap. Program constants, report precision, workspace and physical time are distinct from label count.
